% assets/a3-az-polar.svg — 방향각의 부호 정의와 az_polar 의 네 규칙
%
%   bash _tools/tikz2svg.sh assets/src/a3-az-polar.tex assets/a3-az-polar.svg
%
% 이 그림이 말하는 것
%   (가) 델타를 어디서 어느 쪽으로 재는가 — 부호 정의와 쐐기
%   (나) 쐐기 밖의 각도를 "반대로 틀고 뒤로 민다" 로 바꿔 적는 과정 (규칙 a)
%   (다) az_polar 이 성분 둘을 지령 둘로 바꿀 때 거치는 네 규칙의 순서
%
% 화면 좌표 — (가)(나) 모두 같은 규칙으로 그린다
%   화면 x = 우현 쪽 (y_b 또는 F_y), 화면 y = 선수 쪽 (x_b 또는 F_x)
%   델타는 x_b 에서 y_b 쪽으로 재므로 화면에서는 **위에서 오른쪽으로**(시계방향)
%   화면 극각 theta = 90 - delta  — 아래 모든 arc 가 이 관계를 쓴다
%       delta =   0  ->  theta =  90   (위)
%       delta = +45  ->  theta =  45   (오른위)
%       delta = -45  ->  theta = 135   (왼위)
%   따라서 쐐기는 theta 가 45~135 인 위쪽 부채꼴이다
%
% (나) 의 숫자는 1-7 절 예제 3 의 우현 추진기다
%   (F_x, F_y) = (-50, 100) N  ->  delta_raw = 116.565 deg  (쐐기 밖)
%   접으면 T = -111.803 N, delta = -63.435 deg  — 여전히 쐐기 밖이라 (c) 가 깎는다
%   눈금 0.2 mm/N — 100 N 이 20 mm

\documentclass[border=6pt]{standalone}
% capstone-style.tex — 이 볼트의 TikZ 그림이 공유하는 스타일
%
% 쓰는 법:  % capstone-style.tex — 이 볼트의 TikZ 그림이 공유하는 스타일
%
% 쓰는 법:  % capstone-style.tex — 이 볼트의 TikZ 그림이 공유하는 스타일
%
% 쓰는 법:  \input{capstone-style}   (assets/src/ 안에서)
% 빌드:     bash _tools/tikz2svg.sh assets/src/w02-node-topic.tex assets/w02-node-topic.svg
%
% 왜 TikZ 인가
%   손으로 SVG 좌표를 쓰면 화살촉·선 굵기·글자 크기가 그림마다 미묘하게 달라진다.
%   그것이 "자료가 어설퍼 보인다"의 정체다. 스타일을 한 번만 정하고 상대좌표로
%   배치하면 좌표를 고쳐도 그림이 무너지지 않는다.
%   수식은 본문과 같은 TeX 글꼴로 나온다.
%
% 한글
%   ko.TeX(DVI 경로)를 쓴다. XeLaTeX 은 TikZ 그래픽을 PDF special 로 내보내는데
%   dvisvgm 이 그것을 받으려면 Ghostscript < 10.01 이 필요하다. 이 PC 에는
%   10.07 이 깔려 있어 그 경로가 막힌다. latex -> DVI -> dvisvgm 이 유일하게
%   동작하는 조합이다.

\usepackage[hangul]{kotex}
\usepackage{tikz}
\usepackage{amsmath}
\usepackage{xcolor}
\usetikzlibrary{arrows.meta, positioning, calc, fit, backgrounds, decorations.pathmorphing}

% ---- 색 --------------------------------------------------------------------
% 강조는 하나만. 나머지는 검정. 흑백 인쇄에서도 읽혀야 한다.
\definecolor{ink}{HTML}{1A1A1A}
\definecolor{soft}{HTML}{666666}
\definecolor{accent}{HTML}{7C3AED}
\definecolor{warn}{HTML}{B91C1C}

% 계층 색 — Simulink 모델의 블록 색과 같은 규칙을 쓴다
\definecolor{guid}{HTML}{1D4ED8}    % 유도
\definecolor{ctrl}{HTML}{EA580C}    % 제어
\definecolor{thr}{HTML}{D97706}     % 추진기
\definecolor{plant}{HTML}{16A34A}   % 운동모델
\definecolor{nav}{HTML}{7C3AED}     % 항법 · 신호처리
\definecolor{pathc}{HTML}{0D9488}   % 계획 경로 · 항적
\definecolor{errc}{HTML}{C2410C}    % 오차
\definecolor{flow}{HTML}{0369A1}    % 조류 · 바람

% ---- 선과 화살촉 -----------------------------------------------------------
\tikzset{
  sig/.style   = {ink, line width=0.5pt, -{Stealth[length=4pt, width=3pt]}},
  wire/.style  = {ink, line width=0.5pt},
  fb/.style    = {ink, line width=0.5pt, densely dashed,
                  -{Stealth[length=4pt, width=3pt]}},
  asig/.style  = {accent, line width=0.5pt, densely dashed,
                  -{Stealth[length=4pt, width=3pt]}},
  % 블록 — 흰 바탕, 직각, 검은 테두리. 색은 테두리에만 준다
  blk/.style   = {draw=ink, line width=0.55pt, fill=white,
                  minimum width=18mm, minimum height=8.5mm, inner sep=2pt, font=\small},
  gblk/.style  = {blk, draw=guid},
  cblk/.style  = {blk, draw=ctrl},
  tblk/.style  = {blk, draw=thr},
  pblk/.style  = {blk, draw=plant},
  nblk/.style  = {blk, draw=nav},
  % 합산점 · 분기점
  sum/.style   = {draw=ink, line width=0.55pt, fill=white, circle,
                  inner sep=0pt, minimum size=4.6mm, font=\scriptsize},
  dot/.style   = {ink, fill=ink, circle, inner sep=0pt, minimum size=1.5mm},
  % 글자
  sname/.style = {font=\small, inner sep=1.5pt},
  note/.style  = {font=\scriptsize, soft, inner sep=1.5pt},
  wnote/.style = {font=\scriptsize, warn, inner sep=1.5pt},
  anote/.style = {font=\scriptsize, accent, inner sep=1.5pt},
  axis/.style  = {ink, line width=0.5pt},
  tick/.style  = {font=\scriptsize, soft},
  % 묶음 상자 — 서브시스템 표시
  grp/.style   = {draw=soft, line width=0.4pt, densely dotted, rounded corners=1mm,
                  inner sep=3mm},
}
   (assets/src/ 안에서)
% 빌드:     bash _tools/tikz2svg.sh assets/src/w02-node-topic.tex assets/w02-node-topic.svg
%
% 왜 TikZ 인가
%   손으로 SVG 좌표를 쓰면 화살촉·선 굵기·글자 크기가 그림마다 미묘하게 달라진다.
%   그것이 "자료가 어설퍼 보인다"의 정체다. 스타일을 한 번만 정하고 상대좌표로
%   배치하면 좌표를 고쳐도 그림이 무너지지 않는다.
%   수식은 본문과 같은 TeX 글꼴로 나온다.
%
% 한글
%   ko.TeX(DVI 경로)를 쓴다. XeLaTeX 은 TikZ 그래픽을 PDF special 로 내보내는데
%   dvisvgm 이 그것을 받으려면 Ghostscript < 10.01 이 필요하다. 이 PC 에는
%   10.07 이 깔려 있어 그 경로가 막힌다. latex -> DVI -> dvisvgm 이 유일하게
%   동작하는 조합이다.

\usepackage[hangul]{kotex}
\usepackage{tikz}
\usepackage{amsmath}
\usepackage{xcolor}
\usetikzlibrary{arrows.meta, positioning, calc, fit, backgrounds, decorations.pathmorphing}

% ---- 색 --------------------------------------------------------------------
% 강조는 하나만. 나머지는 검정. 흑백 인쇄에서도 읽혀야 한다.
\definecolor{ink}{HTML}{1A1A1A}
\definecolor{soft}{HTML}{666666}
\definecolor{accent}{HTML}{7C3AED}
\definecolor{warn}{HTML}{B91C1C}

% 계층 색 — Simulink 모델의 블록 색과 같은 규칙을 쓴다
\definecolor{guid}{HTML}{1D4ED8}    % 유도
\definecolor{ctrl}{HTML}{EA580C}    % 제어
\definecolor{thr}{HTML}{D97706}     % 추진기
\definecolor{plant}{HTML}{16A34A}   % 운동모델
\definecolor{nav}{HTML}{7C3AED}     % 항법 · 신호처리
\definecolor{pathc}{HTML}{0D9488}   % 계획 경로 · 항적
\definecolor{errc}{HTML}{C2410C}    % 오차
\definecolor{flow}{HTML}{0369A1}    % 조류 · 바람

% ---- 선과 화살촉 -----------------------------------------------------------
\tikzset{
  sig/.style   = {ink, line width=0.5pt, -{Stealth[length=4pt, width=3pt]}},
  wire/.style  = {ink, line width=0.5pt},
  fb/.style    = {ink, line width=0.5pt, densely dashed,
                  -{Stealth[length=4pt, width=3pt]}},
  asig/.style  = {accent, line width=0.5pt, densely dashed,
                  -{Stealth[length=4pt, width=3pt]}},
  % 블록 — 흰 바탕, 직각, 검은 테두리. 색은 테두리에만 준다
  blk/.style   = {draw=ink, line width=0.55pt, fill=white,
                  minimum width=18mm, minimum height=8.5mm, inner sep=2pt, font=\small},
  gblk/.style  = {blk, draw=guid},
  cblk/.style  = {blk, draw=ctrl},
  tblk/.style  = {blk, draw=thr},
  pblk/.style  = {blk, draw=plant},
  nblk/.style  = {blk, draw=nav},
  % 합산점 · 분기점
  sum/.style   = {draw=ink, line width=0.55pt, fill=white, circle,
                  inner sep=0pt, minimum size=4.6mm, font=\scriptsize},
  dot/.style   = {ink, fill=ink, circle, inner sep=0pt, minimum size=1.5mm},
  % 글자
  sname/.style = {font=\small, inner sep=1.5pt},
  note/.style  = {font=\scriptsize, soft, inner sep=1.5pt},
  wnote/.style = {font=\scriptsize, warn, inner sep=1.5pt},
  anote/.style = {font=\scriptsize, accent, inner sep=1.5pt},
  axis/.style  = {ink, line width=0.5pt},
  tick/.style  = {font=\scriptsize, soft},
  % 묶음 상자 — 서브시스템 표시
  grp/.style   = {draw=soft, line width=0.4pt, densely dotted, rounded corners=1mm,
                  inner sep=3mm},
}
   (assets/src/ 안에서)
% 빌드:     bash _tools/tikz2svg.sh assets/src/w02-node-topic.tex assets/w02-node-topic.svg
%
% 왜 TikZ 인가
%   손으로 SVG 좌표를 쓰면 화살촉·선 굵기·글자 크기가 그림마다 미묘하게 달라진다.
%   그것이 "자료가 어설퍼 보인다"의 정체다. 스타일을 한 번만 정하고 상대좌표로
%   배치하면 좌표를 고쳐도 그림이 무너지지 않는다.
%   수식은 본문과 같은 TeX 글꼴로 나온다.
%
% 한글
%   ko.TeX(DVI 경로)를 쓴다. XeLaTeX 은 TikZ 그래픽을 PDF special 로 내보내는데
%   dvisvgm 이 그것을 받으려면 Ghostscript < 10.01 이 필요하다. 이 PC 에는
%   10.07 이 깔려 있어 그 경로가 막힌다. latex -> DVI -> dvisvgm 이 유일하게
%   동작하는 조합이다.

\usepackage[hangul]{kotex}
\usepackage{tikz}
\usepackage{amsmath}
\usepackage{xcolor}
\usetikzlibrary{arrows.meta, positioning, calc, fit, backgrounds, decorations.pathmorphing}

% ---- 색 --------------------------------------------------------------------
% 강조는 하나만. 나머지는 검정. 흑백 인쇄에서도 읽혀야 한다.
\definecolor{ink}{HTML}{1A1A1A}
\definecolor{soft}{HTML}{666666}
\definecolor{accent}{HTML}{7C3AED}
\definecolor{warn}{HTML}{B91C1C}

% 계층 색 — Simulink 모델의 블록 색과 같은 규칙을 쓴다
\definecolor{guid}{HTML}{1D4ED8}    % 유도
\definecolor{ctrl}{HTML}{EA580C}    % 제어
\definecolor{thr}{HTML}{D97706}     % 추진기
\definecolor{plant}{HTML}{16A34A}   % 운동모델
\definecolor{nav}{HTML}{7C3AED}     % 항법 · 신호처리
\definecolor{pathc}{HTML}{0D9488}   % 계획 경로 · 항적
\definecolor{errc}{HTML}{C2410C}    % 오차
\definecolor{flow}{HTML}{0369A1}    % 조류 · 바람

% ---- 선과 화살촉 -----------------------------------------------------------
\tikzset{
  sig/.style   = {ink, line width=0.5pt, -{Stealth[length=4pt, width=3pt]}},
  wire/.style  = {ink, line width=0.5pt},
  fb/.style    = {ink, line width=0.5pt, densely dashed,
                  -{Stealth[length=4pt, width=3pt]}},
  asig/.style  = {accent, line width=0.5pt, densely dashed,
                  -{Stealth[length=4pt, width=3pt]}},
  % 블록 — 흰 바탕, 직각, 검은 테두리. 색은 테두리에만 준다
  blk/.style   = {draw=ink, line width=0.55pt, fill=white,
                  minimum width=18mm, minimum height=8.5mm, inner sep=2pt, font=\small},
  gblk/.style  = {blk, draw=guid},
  cblk/.style  = {blk, draw=ctrl},
  tblk/.style  = {blk, draw=thr},
  pblk/.style  = {blk, draw=plant},
  nblk/.style  = {blk, draw=nav},
  % 합산점 · 분기점
  sum/.style   = {draw=ink, line width=0.55pt, fill=white, circle,
                  inner sep=0pt, minimum size=4.6mm, font=\scriptsize},
  dot/.style   = {ink, fill=ink, circle, inner sep=0pt, minimum size=1.5mm},
  % 글자
  sname/.style = {font=\small, inner sep=1.5pt},
  note/.style  = {font=\scriptsize, soft, inner sep=1.5pt},
  wnote/.style = {font=\scriptsize, warn, inner sep=1.5pt},
  anote/.style = {font=\scriptsize, accent, inner sep=1.5pt},
  axis/.style  = {ink, line width=0.5pt},
  tick/.style  = {font=\scriptsize, soft},
  % 묶음 상자 — 서브시스템 표시
  grp/.style   = {draw=soft, line width=0.4pt, densely dotted, rounded corners=1mm,
                  inner sep=3mm},
}


\begin{document}
\begin{tikzpicture}[x=1mm, y=1mm]

\tikzset{
  sbx/.style = {rectangle, rounded corners=1pt, draw, line width=0.7pt,
                align=center, font=\scriptsize, inner sep=2.5pt,
                minimum width=42mm, minimum height=24mm},
  ar/.style  = {-{Stealth[length=4pt,width=3pt]}, ink, line width=0.7pt},
  axs/.style = {nav, line width=0.6pt, -{Stealth[length=3.4pt,width=2.4pt]}},
}

% ===========================================================================
% (가) 방향각의 부호 정의
% ===========================================================================
\begin{scope}[shift={(0,0)}]

  \node[font=\small\bfseries, anchor=north west] at (0,46)
       {(가) $\delta$ 는 선수 방향에서 우현 쪽으로 재는 각이다};

  \begin{scope}[shift={(44,0)}]

    % --- 쐐기 (delta = -45 ~ +45  ->  theta = 45 ~ 135) ------------------
    \fill[guid!10] (0,0) -- (45:30mm) arc (45:135:30mm) -- cycle;
    \draw[guid, line width=0.4pt] (45:30mm) arc (45:135:30mm);

    % --- 좌표축 ----------------------------------------------------------
    \draw[axs] (0,-32) -- (0,34);
    \node[font=\scriptsize, nav, anchor=south] at (0,34.5) {$x_b$ (선수)};
    \draw[axs] (-34,0) -- (34,0);
    \node[font=\scriptsize, nav, anchor=west] at (34.5,0) {$y_b$ (우현)};

    % --- 세 기준선 -------------------------------------------------------
    \draw[soft, line width=0.4pt, densely dashed] (0,0) -- (45:32mm);
    \draw[soft, line width=0.4pt, densely dashed] (0,0) -- (135:32mm);
    \node[font=\scriptsize, guid, anchor=south west] at (45:33mm) {$\delta=+45^\circ$};
    \node[font=\scriptsize, guid, anchor=south east] at (135:33mm) {$\delta=-45^\circ$};
    \node[font=\scriptsize, ink, anchor=south west, inner sep=1pt] at (1.5,30)
         {$\delta=0$ — 똑바로 앞으로 밂};

    % --- 추력 벡터 하나 — delta = +30  ->  theta = 60 ---------------------
    \draw[thr, line width=1.4pt, -{Stealth[length=4.4pt,width=3.2pt]}]
         (0,0) -- (60:26mm);
    \node[font=\scriptsize, thr, anchor=north west, inner sep=1pt]
         at (55:17mm) {$T>0$};

    % --- delta 호 — 위(90)에서 벡터(60)까지 ------------------------------
    \draw[ctrl, line width=0.5pt, -{Stealth[length=2.6pt,width=2pt]}]
         (90:13mm) arc (90:60:13mm);
    \node[font=\scriptsize, ctrl, anchor=south west, inner sep=1pt] at (73:14.5mm)
         {$\delta$};

    % --- 같은 각도에서 T < 0 — 반대 방향 ---------------------------------
    \draw[errc, line width=0.9pt, densely dashed,
          -{Stealth[length=4pt,width=3pt]}] (0,0) -- (240:22mm);
    \node[font=\scriptsize, errc, anchor=north east, inner sep=1pt] at (240:23mm)
         {$T<0$ — 같은 $\delta$, 뒤로 밂};

    % --- 추진기 위치 표시 -------------------------------------------------
    \node[circle, draw=thr, fill=white, line width=0.7pt, minimum size=3.4mm,
          inner sep=0pt] at (0,0) {};
  \end{scope}

  \node[anchor=north west, font=\scriptsize, text width=88mm] at (0,-38)
       {\textbf{부호를 외우지 말고 이렇게 읽는다} \\[4pt]
        $\delta$ 는 \textbf{힘이 향하는 방향}이다 — 노즐이 아니라 힘의 방향으로 정했다 \\[2pt]
        \hspace*{4mm}$\delta = 0$ \ 앞으로 밂 (전진) \\[2pt]
        \hspace*{4mm}$\delta > 0$ \ 우현 쪽으로 기울어 밂 \\[2pt]
        \hspace*{4mm}$\delta < 0$ \ 좌현 쪽으로 기울어 밂 \\[4pt]
        $T$ 의 부호는 \textbf{방향을 뒤집는다} — 같은 $\delta$ 에서 $T<0$ 이면
        화살표가 정반대를 가리킨다 \\[4pt]
        $\;\to\;$ 그래서 한 힘을 적는 방법이 \textbf{언제나 둘}이다.
        $(T,\ \delta)$ 와 $(-T,\ \delta \mp 180^\circ)$ 가 같은 힘이다.
        (나)가 그 둘 중 쐐기에 가까운 쪽을 고르는 일이다};
\end{scope}

% ===========================================================================
% (나) 규칙 (a) — 뒤쪽 90 도 접기
% ===========================================================================
\begin{scope}[shift={(98,0)}]

  \node[font=\small\bfseries, anchor=north west] at (0,46)
       {(나) 쐐기 밖의 각도를 "뒤로 밂" 으로 바꿔 적는다};

  \begin{scope}[shift={(44,0)}]

    % --- 앞쪽 쐐기와 뒤쪽 쐐기 -------------------------------------------
    \fill[guid!10] (0,0) -- (45:30mm) arc (45:135:30mm) -- cycle;
    \fill[errc!8]  (0,0) -- (225:30mm) arc (225:315:30mm) -- cycle;
    \draw[guid, line width=0.4pt] (45:30mm) arc (45:135:30mm);
    \draw[errc, line width=0.4pt, densely dashed] (225:30mm) arc (225:315:30mm);
    \node[font=\scriptsize, guid, anchor=south] at (90:31mm) {$T>0$ 로 낼 수 있는 쐐기};
    \node[font=\scriptsize, errc, anchor=north] at (270:31mm) {$T<0$ 로 낼 수 있는 쐐기};

    % --- 좌표축 — F_y 가 화면 오른쪽, F_x 가 화면 위 --------------------
    \draw[axs] (0,-32) -- (0,34);
    \node[font=\scriptsize, nav, anchor=south] at (0,34.5) {$F_x$ [N]};
    \draw[axs] (-34,0) -- (34,0);
    \node[font=\scriptsize, nav, anchor=west] at (34.5,0) {$F_y$ [N]};

    % --- 힘 벡터 — (F_x, F_y) = (-50, 100) -> 화면 (20, -10) ------------
    \draw[thr, line width=1.4pt, -{Stealth[length=4.4pt,width=3.2pt]}]
         (0,0) -- (20,-10);
    \node[font=\scriptsize, thr, anchor=north west, inner sep=1pt] at (20.5,-10.5)
         {$(F_x, F_y) = (-50,\ 100)$};

    % --- 접기 전 각도 호 — 위(90)에서 벡터(-26.565)까지 -----------------
    \draw[warn, line width=0.5pt, -{Stealth[length=2.6pt,width=2pt]}]
         (90:11mm) arc (90:-26.565:11mm);
    \node[font=\scriptsize, warn, anchor=west, inner sep=1pt] at (40:12.5mm)
         {$\delta = 116.565^\circ$};

    % --- 접은 뒤의 노즐 방향 — delta = -63.435 -> theta = 153.435 -------
    \draw[plant, line width=1.1pt, densely dashed,
          -{Stealth[length=4pt,width=3pt]}] (0,0) -- (153.435:24mm);
    \node[font=\scriptsize, plant, anchor=south east, inner sep=1pt]
         at (153.435:25mm) {접은 각도 $\delta = -63.435^\circ$};

    \node[circle, draw=thr, fill=white, line width=0.7pt, minimum size=3.4mm,
          inner sep=0pt] at (0,0) {};
  \end{scope}

  \node[anchor=north west, font=\scriptsize, text width=88mm] at (0,-38)
       {\textbf{접기는 힘을 바꾸지 않는다 — 적는 법만 바꾼다} \\[4pt]
        \hspace*{4mm}접기 전 \quad $T = +111.803$ N, \ $\delta = +116.565^\circ$ \\[2pt]
        \hspace*{4mm}접은 뒤 \quad $T = -111.803$ N, \ $\delta = -63.435^\circ$ \\[4pt]
        두 줄이 가리키는 힘은 \textbf{같은 벡터}다. 서보를 반 바퀴 돌리는 대신
        프로펠러를 거꾸로 돌린다 \\[4pt]
        $-63.435^\circ$ 는 여전히 $\pm45^\circ$ 밖이다
        $\;\to\;$ 규칙 (c)가 $-45^\circ$ 로 깎고, 깎인 만큼이 못 낸 힘이 된다.
        1-7절의 null space 보정이 필요해지는 자리가 바로 여기다};
\end{scope}

% ===========================================================================
% (다) az_polar 의 네 규칙 — 거치는 순서
% ===========================================================================
\draw[wire, soft] (0,-78) -- (186,-78);

\node[font=\small\bfseries, anchor=north west] at (0,-83)
     {(다) \texttt{az\_polar} 이 성분 둘을 지령 둘로 바꾸는 순서};

\node[sbx, draw=guid, fill=guid!6] (p1) at (22,-108)
     {\textbf{(a) 뒤쪽 $90^\circ$ 접기} \\[2pt]
      $\lvert\delta\rvert > 90^\circ$ 이면 \\
      $\delta \mathrel{-}= \operatorname{sign}(\delta)\,\pi$, \ $T \mathrel{\cdot}= -1$};

\node[sbx, draw=soft, fill=white] (p2) at (70,-108)
     {\textbf{(b) 반전 금지} \\[2pt]
      \texttt{rev < 0.5} 이면 \\ (a)를 되돌림};

\node[sbx, draw=warn, fill=warn!6] (p3) at (118,-108)
     {\textbf{(c) 각도 한계} \\[2pt]
      $\delta \leftarrow \operatorname{clip}(\delta, \pm d_{\max})$ \\
      추력은 건드리지 않음};

\node[sbx, draw=warn, fill=warn!6] (p4) at (166,-108)
     {\textbf{(d) 추력 한계} \\[2pt]
      $T \leftarrow \operatorname{clip}(T, \pm F_{\max})$ \\
      비율로 줄이지 않음};

\draw[ar] (p1) -- (p2);
\draw[ar] (p2) -- (p3);
\draw[ar] (p3) -- (p4);

\node[font=\scriptsize, soft, anchor=south west] at (1,-93)
     {입력 $(F_x, F_y)$ \ $\to$ \ $T=\sqrt{F_x^2+F_y^2}$, \
      $\delta=\operatorname{atan2}(F_y,F_x)$};
\node[font=\scriptsize, soft, anchor=south east] at (187,-93) {출력 $(T,\ \delta)$};

\node[anchor=north west, font=\scriptsize, text width=88mm] at (0,-126)
     {\textbf{(a)(b)는 요구를 지킨다} \\[4pt]
      같은 힘을 다르게 적는 것뿐이므로 렌치가 바뀌지 않는다.
      \texttt{rev = 0} (반전 금지)일 때만 (b)가 (a)를 되돌리고,
      그 결과 각도가 한계를 한참 넘어 (c)에서 크게 잘린다};

\node[anchor=north west, font=\scriptsize, text width=88mm] at (98,-126)
     {\textcolor{warn}{\textbf{(c)(d)는 요구를 깎는다 — 조용히}} \\[4pt]
      오류도 경고도 뜨지 않는다. 잘렸다는 사실은 1-10절의 요구 대비 실제
      대조 그림에서만 보인다. 요구 전체를 같은 비율로 줄이는 처리를 두지
      않은 이유는, 어느 축이 왜 모자랐는지가 그림에서 지워지기 때문이다};

\end{tikzpicture}
\end{document}
